\section*{Capítulo \ref{ch:g12:innate-immunity} --- La inmunidad innata}

\begin{solution}{exo:g12:innate-immunity:1}
Innata: presente desde el nacimiento, actúa en unas horas, reconoce
clases de intrusos por patrones compartidos, responde igual todas las
veces. Adaptativa: tarda días, reconoce a cada intruso individualmente,
lo recuerda (y solo existe en los \omterm{def:g10:common-ancestry:bodyplan}{vertebrados}).
\end{solution}

\begin{solution}{exo:g12:innate-immunity:2}
Enrojecimiento y calor: unos vasos ensanchados que traen más sangre.
Hinchazón: unas paredes vasculares que dejan escapar plasma al tejido.
Dolor: unas terminaciones nerviosas sensibilizadas por los \omterm{prop:g12:innate-immunity:inflammation}{mediadores}.
\end{solution}

\begin{solution}{exo:g12:innate-immunity:3}
Patrones moleculares compartidos por clases de microbios (fragmentos de
pared, \omterm{def:g10:chemistry-of-life:families}{ácidos nucleicos} víricos) y moléculas liberadas por las \omterm{def:g10:cells-common-unit:cell}{células}
dañadas. Liberan \omterm{prop:g12:innate-immunity:inflammation}{mediadores}: histamina, prostaglandinas, citocinas.
\end{solution}

\begin{solution}{exo:g12:innate-immunity:4}
Contacto y fijación por los receptores; englobamiento por la membrana;
digestión en una vesícula con \omterm{def:g11:enzymes-and-phenotype:enzyme}{enzimas} y sustancias químicas reactivas;
exhibición de fragmentos proteicos en la superficie.
\end{solution}

\begin{solution}{exo:g12:innate-immunity:5}
Bloquean la \omterm{def:g11:enzymes-and-phenotype:enzyme}{enzima} que fabrica las prostaglandinas, con lo que reducen
la hinchazón, el dolor y la fiebre. No suprimen la causa ni matan a
ningún microbio.
\end{solution}

\begin{solution}{exo:g12:innate-immunity:6}
La innata, hacia el medio día; la adaptativa, hacia el día 9; las
curvas se cruzan alrededor de los días 5--6.
\end{solution}

\begin{solution}{exo:g12:innate-immunity:7}
El disparador es un patrón molecular reconocido por las \omterm{def:g10:cells-common-unit:cell}{células}
centinela, no un intruso vivo. El suero salino estéril no lleva ningún
patrón ni causa daño: como mucho, una respuesta ligera a la aguja.
\end{solution}

\begin{solution}{exo:g12:innate-immunity:8}
Neutrófilos muertos, bacterias muertas y vivas, restos digeridos y
plasma. Los neutrófilos llegan a lo largo de unas horas y mueren a lo
largo de un día; el pus es lo que se acumula una vez que han peleado.
\end{solution}

\begin{solution}{exo:g12:innate-immunity:9}
La fiebre del heno es histamina liberada por los mastocitos contra el
polen: bloquear la histamina suprime la respuesta. En una infección, la
histamina impulsa solo los primeros minutos; la hinchazón, el dolor y
el reclutamiento posteriores funcionan con prostaglandinas y citocinas,
que el antihistamínico no toca.
\end{solution}

\begin{solution}{exo:g12:innate-immunity:10}
A las 12 horas, 100\% frente a 55\%; a las 36 horas, 50\% frente a
30\%. El medicamento baja la hinchazón en todo momento, pero la curva
sigue cayendo a cero hacia las 72 horas: no acorta la respuesta.
\end{solution}

\begin{solution}{exo:g12:innate-immunity:11}
Las fibras desgarradas del \omterm{def:g10:muscles-and-joints:joint}{ligamento} y las \omterm{def:g10:cells-common-unit:cell}{células} aplastadas liberan
las moléculas propias de las \omterm{def:g10:cells-common-unit:cell}{células} dañadas, que los receptores de las
\omterm{def:g10:cells-common-unit:cell}{células} centinela reconocen igual que reconocerían a un microbio: el
daño por sí solo es un disparador.
\end{solution}

\begin{solution}{exo:g12:innate-immunity:12}
Sin neutrófilos se pierden los primeros días y unas bacterias
corrientes desbordan al organismo: la respuesta innata es la que
aguanta la línea de inmediato. Sin \omterm{def:g12:innate-immunity:immunity}{inmunidad adaptativa}, la respuesta
innata despacha las primeras infecciones, pero las persistentes o
repetidas, que necesitan anticuerpos específicos y memoria, acaban
ganando: la respuesta adaptativa remata lo que la innata contiene.
\end{solution}

\begin{solution}{exo:g12:innate-immunity:13}
Un absceso es una inflamación que pelea contra bacterias vivas:
amortiguarla frena su eliminación y arriesga que se extiendan. La
artritis es una inflamación sin ningún microbio, dirigida contra la
propia \omterm{def:g10:muscles-and-joints:joint}{articulación}: ahí la respuesta es la enfermedad, y reducirla es
el tratamiento.
\end{solution}

\begin{solution}{exo:g12:innate-immunity:14}
Los receptores innatos los fijan los \omterm{def:g10:universal-dna:gene}{genes} y reconocen patrones
compartidos por clases enteras; nada de un segundo encuentro los
cambia. El sistema adaptativo construye receptores específicos para
cada intruso y conserva las \omterm{def:g10:cells-common-unit:cell}{células} que los fabricaron: reconoce al
individuo, y conserva el reconocimiento.
\end{solution}

\begin{solution}{exo:g12:innate-immunity:15}
La inflamación es la cura cuando va dirigida contra un intruso o una
lesión reales, es proporcionada y se resuelve en unos días: limpia el
sitio y pone en marcha la reparación. Es la enfermedad cuando va
dirigida contra algo inofensivo (el polen, las propias \omterm{def:g10:muscles-and-joints:joint}{articulaciones}
del cuerpo), es excesiva o no termina: entonces sus \omterm{prop:g12:innate-immunity:inflammation}{mediadores} y sus
\omterm{def:g12:innate-immunity:phagocytosis}{fagocitos} dañan el tejido que debían proteger. Lo que decide es la
diana y la duración, no el mecanismo, que es el mismo en los dos casos.
\end{solution}

\begin{solution}{pb:g12:innate-immunity:1}
\textbf{1.} Tres duplicaciones en 2 horas: $\num{10000} \times 8 =
\num{80000}$; seis en 4 horas: $\num{640000}$.

\textbf{2.} $200 \times 5 = 1000$ por hora --- muy por debajo de las
$\num{10000}$ bacterias nuevas que añaden las duplicaciones de la
primera hora: el crecimiento continúa.

\textbf{3.} $\num{50000} \times 20 = 10^6$ bacterias.

\textbf{4.} En la hora 2 hay menos de $\num{80000}$ bacterias; los
neutrófilos de la primera hora pueden retirar por sí solos un millón.
Los \omterm{def:g12:innate-immunity:phagocytosis}{fagocitos} van ganando a partir de la hora 2.

\textbf{5.} En unas pocas horas las bacterias han desaparecido; lo que
queda son los neutrófilos muertos, los restos y el líquido que serán el
pus, y los macrófagos que lo retiran.

\textbf{6.} Enrojecimiento y calor: la histamina y las prostaglandinas
ensanchan los vasos. Hinchazón: la histamina hace que las paredes dejen
escapar líquido y el plasma se filtra. Dolor: las prostaglandinas
sensibilizan las terminaciones nerviosas.

\textbf{7.} Esos vasos que dejan escapar líquido son la vía por la que
los neutrófilos y las \omterm{def:g10:chemistry-of-life:families}{proteínas} defensivas del plasma llegan al tejido;
el líquido que aprieta los nervios sensibilizados es el dolor.

\textbf{8.} Cada latido empuja más sangre a los vasos ensanchados; el
tejido hinchado no puede expandirse, así que cada pulso sube la presión
sobre los nervios sensibilizados: una palpitación al compás del
corazón.

\textbf{9.} Neutrófilos muertos, bacterias muertas y vivas, restos
digeridos, plasma.

\textbf{10.} Los patrones desaparecieron con las bacterias y los restos
se retiraron: sin patrón no hay reconocimiento por las centinelas, y
sin él no hay \omterm{prop:g12:innate-immunity:inflammation}{mediadores}; los vasos vuelven a la normalidad.

\textbf{11.} Reduce el enrojecimiento, el calor y la fuga precoces
debidos a la histamina; no la hinchazón ni el dolor posteriores,
debidos a las prostaglandinas, ni el reclutamiento de neutrófilos.

\textbf{12.} Bloquea la \omterm{def:g11:enzymes-and-phenotype:enzyme}{enzima} que fabrica las prostaglandinas: bajan
el dolor, parte de la hinchazón y la fiebre.

\textbf{13.} La crema apaga casi todos los \omterm{prop:g12:innate-immunity:inflammation}{mediadores}, así que se llama
a menos neutrófilos; con menos \omterm{def:g12:innate-immunity:phagocytosis}{fagocitos}, puede que las bacterias no se
eliminen y que la infección se extienda.

\textbf{14.} Los tres actúan sobre los signos y ninguno mata bacterias;
el corticoide, al recortar el reclutamiento de \omterm{def:g12:innate-immunity:phagocytosis}{fagocitos}, puede hacer
que la infección dure más.

\textbf{15.} Unas citocinas que llegan al cerebro y suben la consigna
de temperatura; el calor de más acelera a los \omterm{def:g12:innate-immunity:phagocytosis}{fagocitos} y frena a las
bacterias.

\textbf{16.} Las citocinas que captaron en el sitio, que les dicen a
los linfocitos qué clase de intruso era; informan a los linfocitos del
ganglio linfático.

\textbf{17.} Alrededor de una semana o diez días. Para una astilla que
la respuesta innata elimina en unas horas no harán falta --- pero se
fabricarán, y se recordarán.

\textbf{18.} La respuesta adaptativa es más rápida y más fuerte la
segunda vez, gracias a la memoria; la respuesta innata es exactamente
la misma.

\textbf{19.} Los \omterm{def:g11:antibiotic-resistance:antibiotic}{antibióticos} matan bacterias y ocupan el lugar de los
\omterm{def:g12:innate-immunity:phagocytosis}{fagocitos} que faltan frente a las infecciones bacterianas; no pueden
sustituir la retirada de restos que hacen los \omterm{def:g12:innate-immunity:phagocytosis}{fagocitos}, ni su acción
sobre los hongos, ni el informe al ganglio linfático.

\textbf{20.} Hacia la hora 2 los \omterm{def:g12:innate-immunity:phagocytosis}{fagocitos} adelantan a las bacterias;
la histamina y las prostaglandinas detrás del enrojecimiento, el calor
y la hinchazón, y las prostaglandinas detrás del dolor; la \omterm{prop:g12:innate-immunity:handover}{célula
dendrítica} es la que lleva el informe.
\end{solution}
